\begin{thesisabstract}
\phantomsection\label{front:abstract}
\abstractparagraph{Seismic performance of hospital buildings depends on the interaction of the superstructure with the supporting soil and foundation. This study evaluates the influence of soil--structure interaction on the G+4 Infectious and Communicable Disease Hospital in Lekhnath, Pokhara, using six reinforced-concrete structural blocks. Fixed-base models were compared with flexible-base representations for hard, medium and soft soil conditions. The response measures considered were fundamental period, modal mass participation, response-spectrum storey displacement, inter-storey drift, base shear, storey shear, foundation movement, vertical reaction and beam and column force demands. X-direction fragility curves were also examined to relate the structural-response findings to seismic vulnerability.}

\abstractparagraph{The fundamental periods increased between fixed-base and soft-soil support by 8.53\%, 19.73\%, 5.22\%, 8.99\%, 9.00\% and 7.85\% for Blocks 1--6, respectively. At least 90\% cumulative modal mass participation in both horizontal directions was reached between Modes 5 and 10 across the investigated cases. The X-direction uppermost-storey response-spectrum displacement increased under softer support in all six blocks; however, the magnitude of amplification differed with structural configuration. The largest reported maximum inter-storey drift was 0.774\% in Block 3 under soft-soil support. Base-shear and storey-shear demands showed smaller or non-monotonic variations for several cases, demonstrating that period lengthening alone does not determine the corresponding seismic force response. The foundation settlements, vertical reactions and structural member-force envelopes associated with the $1.2DL+1.5LL$ load combination also varied among the blocks and support conditions.}

\abstractparagraph{The results show that foundation flexibility can substantially affect deformation-sensitive response measures while producing different changes in force demand from one structural block to another. Consequently, fixed-base assumptions may not fully represent the range of response relevant to hospital performance assessment. The fragility relationships provide supplementary evidence of the sensitivity of damage-state exceedance patterns to foundation-support assumptions. Site-specific validation of soil and foundation properties, producing numerical models and fragility inputs is required before these comparisons are extended to predictions of as-built hospital safety.}

\abstractkeywords{Soil--Structure Interaction; Hospital Buildings; Seismic Response; Foundation Flexibility; Response-Spectrum Analysis; Seismic Fragility.}
\end{thesisabstract}
